\input{configTD}

\title{Analyse Numérique TD4 - Erreurs, EDO et méthode de Newton}

\author{}
\date{}

\newenvironment{solution}
    {\par\vspace{0.5em}\begin{mdframed}[linewidth=0.5pt]\par}
    {\end{mdframed}\par\vspace{0.5em}}

\begin{document}
\sffamily
\maketitle

\section*{Exercice 1 - Euler explicite sur un modèle de refroidissement}
On considère le problème de Cauchy
\[
\begin{cases}
y'(t)=-2y(t),\\
y(0)=1.
\end{cases}
\]
La solution exacte est $y(t)=e^{-2t}$.

\begin{enumerate}
    \item Écrire la formule d'Euler explicite avec un pas $h>0$.

    \item Calculer les approximations $y_1$ et $y_2$ pour $h=0.25$.

    \item Donner la valeur exacte de $y(0.5)$ puis l'erreur absolue commise.

    \item Que devient l'algorithme si on choisit $h=1$ ? Commenter.

    \ifthenelse{\boolean{showSolutions}}{
        \begin{solution}
            Euler explicite donne
            \[
            y_{n+1}=y_n+h(-2y_n)=(1-2h)y_n.
            \]
            Pour $h=0.25$,
            \[
            y_1=(1-0.5)y_0=0.5,
            \qquad
            y_2=(1-0.5)y_1=0.25.
            \]
            La valeur exacte est
            \[
            y(0.5)=e^{-1}\approx 0.3679.
            \]
            L'erreur absolue vaut donc
            \[
            |0.3679-0.25|\approx 0.1179.
            \]
            Si $h=1$, alors
            \[
            y_{n+1}=-y_n.
            \]
            La suite oscille et ne reproduit pas le comportement de décroissance de la solution exacte. Le pas est trop grand.
        \end{solution}
    }{}
\end{enumerate}

\section*{Exercice 2 - Comparer Euler explicite et implicite}
On étudie à nouveau
\[
y'(t)=-2y(t),\qquad y(0)=1.
\]

\begin{enumerate}
    \item Écrire la formule d'Euler implicite.

    \item Calculer $y_1$ et $y_2$ pour $h=1$.

    \item Comparer qualitativement avec Euler explicite sur le meme pas.

    \ifthenelse{\boolean{showSolutions}}{
        \begin{solution}
            Euler implicite donne
            \[
            y_{n+1}=y_n-2hy_{n+1}.
            \]
            Donc
            \[
            y_{n+1}=\frac{1}{1+2h}y_n.
            \]
            Pour $h=1$,
            \[
            y_1=\frac13,\qquad y_2=\frac19.
            \]
            Contrairement à l'explicite, l'implicite garde ici un comportement stable et décroissant même avec un grand pas.
        \end{solution}
    }{}
\end{enumerate}

\section*{Exercice 3 - Différences finies et erreur locale}
On veut approcher la dérivée seconde d'une fonction régulière par
\[
\frac{u(x+h)-2u(x)+u(x-h)}{h^2}.
\]

\begin{enumerate}
    \item À l'aide d'un développement de Taylor de $u(x+h)$ et $u(x-h)$, montrer que
    \[
    \frac{u(x+h)-2u(x)+u(x-h)}{h^2}=u''(x)+O(h^2).
    \]

    \item Interpréter cet ordre d'erreur dans le cadre d'une discrétisation d'une poutre ou d'une paroi.

    \ifthenelse{\boolean{showSolutions}}{
        \begin{solution}
            On écrit
            \[
            u(x+h)=u(x)+hu'(x)+\frac{h^2}{2}u''(x)+\frac{h^3}{6}u^{(3)}(x)+\frac{h^4}{24}u^{(4)}(\xi_1),
            \]
            \[
            u(x-h)=u(x)-hu'(x)+\frac{h^2}{2}u''(x)-\frac{h^3}{6}u^{(3)}(x)+\frac{h^4}{24}u^{(4)}(\xi_2).
            \]
            En additionnant,
            \[
            u(x+h)-2u(x)+u(x-h)=h^2u''(x)+\frac{h^4}{24}\big(u^{(4)}(\xi_1)+u^{(4)}(\xi_2)\big).
            \]
            En divisant par $h^2$,
            \[
            \frac{u(x+h)-2u(x)+u(x-h)}{h^2}=u''(x)+O(h^2).
            \]
            Cela signifie que si l'on divise le pas par $2$, l'erreur locale est en principe divisée par environ $4$.
        \end{solution}
    }{}
\end{enumerate}

\section*{Exercice 4 - Méthode de Newton pour une équation non linéaire}
On cherche un déplacement d'équilibre $u$ vérifiant
\[
f(u)=u+\frac{u^3}{10}-1=0.
\]
On peut voir cette équation comme un modèle simplifié d'équilibre non linéaire.

\begin{enumerate}
    \item Calculer $f'(u)$.

    \item Écrire la formule de Newton associée.

    \item À partir de $u_0=1$, calculer $u_1$ puis $u_2$.

    \item Vérifier numériquement que l'approximation s'améliore rapidement.

    \ifthenelse{\boolean{showSolutions}}{
        \begin{solution}
            On a
            \[
            f'(u)=1+\frac{3u^2}{10}.
            \]
            La formule de Newton est
            \[
            u_{n+1}=u_n-\frac{u_n+\frac{u_n^3}{10}-1}{1+\frac{3u_n^2}{10}}.
            \]
            Avec $u_0=1$,
            \[
            u_1=1-\frac{1+0.1-1}{1+0.3}=1-\frac{0.1}{1.3}=\frac{12}{13}\approx 0.9231.
            \]
            Puis
            \[
            u_2\approx 0.9217.
            \]
            On verifie
            \[
            f(1)=0.1,\qquad
            f(u_1)\approx 0.0017,\qquad
            f(u_2)\approx 10^{-6}.
            \]
            La méthode converge très vite.
        \end{solution}
    }{}
\end{enumerate}

\section*{Exercice 5 - Pourquoi aller plus loin ?}
Dans un problème d'identification de paramètres, on cherche parfois à minimiser une fonction de coût
\[
J(\theta)=\frac12\sum_{i=1}^m \big(y_i^{\mathrm{mes}}-y_i^{\mathrm{mod}}(\theta)\big)^2.
\]

\begin{enumerate}
    \item Expliquer pourquoi cette écriture fait penser aux moindres carrés.

    \item Quel est l'intérêt de Newton ou des méthodes quasi-Newton lorsque $\theta$ contient plusieurs inconnues ?

    \item Dans quel cas une méthode de type Levenberg-Marquardt peut-elle être pertinente ?

    \ifthenelse{\boolean{showSolutions}}{
        \begin{solution}
            La fonction $J$ est la somme des carrés des résidus entre mesures et modèle : c'est exactement l'idée des moindres carrés.
            En dimension plus grande, Newton et les méthodes quasi-Newton exploitent l'information de pente et de courbure pour accélérer la convergence par rapport à une simple descente de gradient.
            Une méthode de type Levenberg-Marquardt est pertinente pour l'ajustement de paramètres dans des problèmes de moindres carrés non linéaires, quand on veut un comportement plus robuste qu'un Newton pur.
        \end{solution}
    }{}
\end{enumerate}

\end{document}
